% ============================================
% MEM-12
% Gerado de molambudos.md — build_miolo.py
% ============================================
\fragtitle{mem}{M}{MEM}{12}{A Rasga Mortalha}
\markboth{MEM-12 --- A Rasga Mortalha}{}

\noindent \textit{O arauto no pátio}

\noindent \textbf{1945 --- Hospital Colônia de Barbacena}

\vspace{0.5em}\noindent\rule{\textwidth}{0.5pt}\vspace{0.5em}

\noindent Vinte e oito anos dentro do Colônia.

\noindent Vinte e oito anos de paredes que descascam, de gritos que viram rotina, de cheiro de creolina que não sai do nariz. Vinte e oito anos sem sentir o vento no rosto sem ser através de uma janela gradeada. Vinte e oito anos de solidão que não cura --- que petrifica.

\noindent Em 1945, eu estava na ala dos silenciados. Uma ala onde ninguém fala porque todos aprenderam que falar não adianta. Eu ficava encolhido num canto, com o diário no colo, escrevendo sem parar. A criatura dentro de mim estava quieta --- hibernando. Eu sentia o peso dela, o calor, a respiração lenta.

\noindent De vez em quando ela se virava dentro do meu peito e eu sentia uma pontada. A pontada não doía. Era \textit{lembrança}. Ela me lembrava que estava lá.

\vspace{0.5em}\noindent\rule{\textwidth}{0.5pt}\vspace{0.5em}

\noindent A Rasga Mortalha apareceu num inverno.

\noindent Junho de 1945. O frio entrava pelas frestas da parede e não saía. Eu dormia com as mãos enfiadas nas axilas, tremendo. Os outros pacientes tremiam também. A gente tremia junto, como bichos numa cerca, esperando o abate.

\noindent Numa noite, um som me acordou.

\noindent Não era grito.
Não era choro.
Era um canto.

\noindent Fino. Agudo. Um canto que subia e descia como uma mulher chorando muito longe.

\noindent \textit{\textquotedblleft{}Rasga mortalha...\textquotedblright{}}, alguém sussurrou ao meu lado.

\noindent Era Seu Nonô. Um velho que estava no Colônia desde 1920 --- 25 anos internado. Ele apontou para a janela com um dedo trêmulo.

\noindent \textit{\textquotedblleft{}Ela chegou.\textquotedblright{}}

\noindent \textit{\textquotedblleft{}Quem?\textquotedblright{}}

\noindent Ele não respondeu. Só apontou.

\noindent No batente da janela, havia uma coruja.

\noindent Ela era branca. Não branca como neve --- branca como osso velho. As penas pareciam luz na escuridão, como se o corpo dela fosse feito de luar condensado. Os olhos eram pretos, redondos, sem pálpebras. Ela me olhou. Virou a cabeça 180 graus --- um movimento lento, fluido, que nenhum pássaro devia ser capaz de fazer.

\noindent \textit{\textquotedblleft{}Rasga mortalha\textquotedblright{}\textit{, repetiu Seu Nonô. }\textquotedblleft{}Quando ela canta, alguém morre. Ela rasga a mortalha do finado com o bico pra alma poder sair. Por isso o nome.\textquotedblright{}}

\noindent A coruja abriu o bico.

\noindent O som que saiu não era um piado.
Era um grito humano.
Um grito de mulher --- mas não de dor.
De \textit{aviso}.

\noindent O grito atravessou a noite. Atravessou as grades. Atravessou o meu peito.

\noindent A criatura dentro de mim \textit{acordou}.

\noindent Eu a senti erguer a cabeça dentro do meu corpo. Senti os olhos dela se abrirem --- os olhos amarelos. Ela \textit{viu} a Rasga Mortalha.

\noindent E a Rasga Mortalha viu ela.

\noindent A coruja inclinou a cabeça. O grito parou. Ficou um silêncio tão absoluto que eu ouvia meu próprio sangue circulando.

\noindent E a Rasga Mortalha falou.

\noindent Não com a voz dela.

\noindent Com a voz da criatura.

\noindent A voz que saía do buraco dentro de mim.

\noindent \textit{\textquotedblleft{}Ainda não.\textquotedblright{}}

\noindent A coruja abriu as asas. Eram enormes. Brancas. Pareciam um lençol flutuando no escuro.

\noindent \textit{\textquotedblleft{}Mas em breve.\textquotedblright{}}

\noindent Ela voou.

\noindent O branco dela sumiu na noite.

\vspace{0.5em}\noindent\rule{\textwidth}{0.5pt}\vspace{0.5em}

\noindent No dia seguinte, Seu Nonô foi encontrado morto na cama.

\noindent O rosto estava virado para a janela.
A boca aberta.
Os olhos abertos.

\noindent As pálpebras tinham sido arrancadas.

\noindent Não sobrou sangue. Não sobrou marca. As pálpebras simplesmente tinham \textit{sumido}. Como se alguém as tivesse levado com uma precisão cirúrgica.

\noindent Os enfermeiros disseram que foram os ratos.

\noindent Eu sabia que não.

\noindent A Rasga Mortalha tinha levado as pálpebras do Seu Nonô.
Para que ele \textit{nunca mais fechasse os olhos}.
Para que ele visse o que estava chegando.
Mesmo depois de morto.

\vspace{0.5em}\noindent\rule{\textwidth}{0.5pt}\vspace{0.5em}

\noindent Na semana seguinte, a coruja voltou.

\noindent Toda noite, ela pousava no batente da janela e me olhava. Não cantava. Não gritava. Ficava ali, imóvel, com os olhos pretos fixos em mim.

\noindent E toda noite, a criatura dentro de mim falava.

\noindent \textit{\textquotedblleft{}Ela é minha língua\textquotedblright{}\textit{, disse. }\textquotedblleft{}Meu olho voador. Minha voz no mundo dos vivos.\textquotedblright{}}

\noindent \textit{\textquotedblleft{}Por que ela veio?\textquotedblright{}}

\noindent \textit{\textquotedblleft{}Pra ver como você está.\textquotedblright{}}

\noindent \textit{\textquotedblleft{}E como eu estou?\textquotedblright{}}

\noindent Silêncio.

\noindent \textit{\textquotedblleft{}Pronto.\textquotedblright{}}

\noindent \textit{\textquotedblleft{}Pronto pra quê?\textquotedblright{}}

\noindent \textit{\textquotedblleft{}Pra passar o diário adiante.\textquotedblright{}}

\vspace{0.5em}\noindent\rule{\textwidth}{0.5pt}\vspace{0.5em}

\noindent Depois daquela noite, eu comecei a escrever de um jeito diferente.

\noindent A letra mudou. As palavras vinham mais rápido. Eu não precisava pensar --- a mão se movia sozinha, como se alguém estivesse guiando a caneta.

\noindent O diário não era mais meu.

\noindent A criatura estava escrevendo através de mim.

\noindent Eu era só a mão.
Ela era a voz.

\noindent E o que ela escrevia --- eu não entendia na hora. Só entendi depois. Quando li de novo. Quando vi que as palavras não eram para mim.

\noindent Eram para quem lesse depois.

\noindent Eram para você.

\vspace{0.5em}\noindent\rule{\textwidth}{0.5pt}\vspace{0.5em}

\noindent Na última noite, a coruja não veio.

\noindent Eu acordei sozinho, na cela escura, e soube que ela tinha ido embora.

\noindent Mas a criatura dentro de mim falou:

\noindent \textit{\textquotedblleft{}Ela foi preparar o caminho.\textquotedblright{}}

\noindent \textit{\textquotedblleft{}Preparar como?\textquotedblright{}}

\noindent \textit{\textquotedblleft{}O diário não pode ficar aqui. Tem que sair. Tem que ser encontrado. Tem que ser aberto.\textquotedblright{}}

\noindent \textit{\textquotedblleft{}Por quem?\textquotedblright{}}

\noindent A criatura riu. A risada ecoou dentro do meu peito.

\noindent \textit{\textquotedblleft{}Alguém que precisa tanto de mim quanto você precisou.\textquotedblright{}}

\noindent \textit{\textquotedblleft{}E quando isso vai acontecer?\textquotedblright{}}

\noindent \textit{\textquotedblleft{}Quando você morrer.\textquotedblright{}}

\noindent \textit{\textquotedblleft{}E até lá?\textquotedblright{}}

\noindent \textit{\textquotedblleft{}A gente espera.\textquotedblright{}}

\noindent Ela ficou quieta.

\noindent \textit{\textquotedblleft{}A gente sempre espera.\textquotedblright{}}

\vspace{0.5em}\noindent\rule{\textwidth}{0.5pt}\vspace{0.5em}

\noindent Foram mais 34 anos de espera.

\noindent 34 anos ouvindo a Rasga Mortalha nos telhados do Colônia.

\noindent 34 anos escrevendo palavras que não eram minhas.

\noindent 34 anos sendo a mão de um monstro.

\noindent Até que, em 1979, o médico novo chegou.

\noindent E o ciclo começou a se fechar.

\vspace{0.5em}
\noindent\textit{\textbf{↪ Rotas:} \rota{MEM-11} \rotasep \rota{MEM-13} \rotasep \rota{DOC-07}}
